\input{../tex-inputs/latex-leseansicht-vorspann}

\section[Gustav Schwarzkopf an Arthur Schnitzler, 9. 8. 1916]{L04251 Gustav Schwarzkopf an Arthur Schnitzler, 9. 8. 1916}
\nopagebreak\mylabel{L04251v}
\rehead{ }\normalsize\beginnumbering\briefempfaengerindex{Schnitzler, Arthur@\textsc{Schnitzler, Arthur}!zzzSchwarzkopf, Gustav@\emph{von Gustav Schwarzkopf}!1916-08-091@{9. 8. 1916}|(be}
\toendnotes[C]{\smallbreak\pagebreak[2]}
\correspDesc{Versand  durch Gustav Schwarzkopf am 9. 8. 1916 in Wien
\newline{}Übermittlung  am 9. 8. 1916 in Wien
\newline{}Erhalt  durch Arthur Schnitzler im Zeitraum [10. 8. 1916
                  – 14. 8. 1916?] in Altaussee}\toendnotes[C]{\smallbreak}
\Standort{CUL, Schnitzler, B 96a.}
\physDesc{Postkarte, 611 Zeichen
\newline{}Handschrift: schwarze Tinte, deutsche Kurrent
\newline{}Versand: Stempel: »\nobreak{}\oindex{I., Innere Stadt@\textbf{I., Innere Stadt}, \emph{Verwaltungsgebiet}|pwk}1/1 Wien 8, 9. VIII. 16, 7\nobreak{}«.  }\toendnotes[C]{\smallbreak}\pstart{}{\pb}I Tiefer Graben 17\oindex{Wien@\textbf{Wien}!I., Innere Stadt@\textbf{I., Innere Stadt}!Tiefer Graben 17@\textbf{Tiefer Graben 17}, \emph{Wohngebäude}|pw}\pend{}\pstart{}II. 6\pend{}{\bigskip}\pstart{}Herrn\pend{}\pstart{}\textsc{D\textsuperscript{r} Arthur Schnitzler}\pend{}\pstart{}\textsc{Alt-Aussee\oindex{Altaussee@\textbf{Altaussee}, \emph{Verwaltungsgebiet}|pw}}\pend{}\pstart{}Oberſteiermark\oindex{Obersteiermark@\textbf{Obersteiermark}|pw}.\pend{}{\bigskip}\vspace{1em}
\pstart
           \raggedleft{}{\pb}\textsc{Wien\oindex{Wien@\textbf{Wien}, \emph{Verwaltungsgebiet}|pw}}{ }9. VIII 16\pend
           
\pstart{}Lieber Arthur,\pend\vspace{0.5em}
\pstart
           da ich gar nichts mehr von Ihnen gehört habe, muß ich annehmen, daß Sie meine \label{K_L04251-1v}\edtext{Karte}{\lemma{\textnormal{\emph{Karte}}}\Cendnote{\textnormal{XXXX Auszeichnungsfehler: Dokument L04252 nicht gefunden.}}}\label{K_L04251-1} aus
                  Marienbad\oindex{Marienbad@\textbf{Marienbad}|pw} nicht erhalten haben. Seit
                  Montag bin ich wieder in Wien\oindex{Wien@\textbf{Wien}, \emph{Verwaltungsgebiet}|pw}; ich
               hatte ja zuerst die kühne Abſicht von M.\oindex{Marienbad@\textbf{Marienbad}|pw} noch
               nach \textsc{Ischl\oindex{Bad Ischl@\textbf{Bad Ischl}|pw}} zu fahren, aber die Schwierigkeiten der Reiſe – man muß jetzt in Linz\oindex{Linz@\textbf{Linz}|pw} übernachten – und die Erfahrungen in M.\oindex{Marienbad@\textbf{Marienbad}|pw} haben mich abgehalten.\pend
           
\pstart
           – Wie geht es ihnen? Hat \textsc{Aussee\oindex{Bad Aussee@\textbf{Bad Aussee}, \emph{Hauptstadt}|pw}} auch ein\textcolor{gray}{en}{ }\label{K_L04251-2v}\edtext{Zauner\oindex{Konditorei Zauner@\textbf{Konditorei Zauner}, \emph{Kaffeehaus}|pw}}{\lemma{\textnormal{\emph{Zauner}}}\Cendnote{\textnormal{Die Konditorei Zauner\oindex{Konditorei Zauner@\textbf{Konditorei Zauner}, \emph{Kaffeehaus}|pwk} gab (und gibt) es in Bad Ischl\oindex{Bad Ischl@\textbf{Bad Ischl}|pwk}, aber nicht in Aussee\oindex{Bad Aussee@\textbf{Bad Aussee}, \emph{Hauptstadt}|pwk}.}}}\label{K_L04251-2}? In M.\oindex{Marienbad@\textbf{Marienbad}|pw} iſt ein
               Zuckerbäcker, den man in reſpektvoller Entfernung ne{\geminationn}en
               darf.\pend
           
\pstart
           {\pb}Ihre Frau\pwindex{Schnitzler, Olga 17.\,1.\,1882 Wien – 13.\,1.\,1970 Lugano@\textsc{Schnitzler, Olga} (17.\,1.\,1882 Wien – 13.\,1.\,1970 Lugano), \emph{Schauspielerin, Sängerin}|pw}, Sie und die{\\[\baselineskip]}Kinder\pwindex{Cappellini, Lili 13.\,9.\,1909 Wien – 26.\,7.\,1928 Venedig@\textsc{Cappellini, Lili} (13.\,9.\,1909 Wien – 26.\,7.\,1928 Venedig)|pwv}\pwindex{Schnitzler, Heinrich 9.\,8.\,1902 Hinterbrühl – 12.\,7.\,1982 Wien@\textsc{Schnitzler, Heinrich} (9.\,8.\,1902 Hinterbrühl – 12.\,7.\,1982 Wien), \emph{Regisseur, Schauspieler}|pwv} herzlichſt{\\[\baselineskip]}grüßend{\\[\baselineskip]}Ihr{\\[\baselineskip]}\spacefill\mbox{Gustav.}\pend
           \leftskip=0em{}\selectlanguage{ngerman}\endnumbering\briefempfaengerindex{Schnitzler, Arthur@\textsc{Schnitzler, Arthur}!zzzSchwarzkopf, Gustav@\emph{von Gustav Schwarzkopf}!1916-08-091@{9. 8. 1916}|)be}\mylabel{L04251h}
\begin{anhang}
\end{anhang}\newcommand{\dateiname}{L04251}\newcommand{\titel}{Gustav Schwarzkopf an Arthur Schnitzler, 9. 8. 1916}\newcommand{\editorInnen}{Herausgegeben von Jahnke, SelmaMüller, Martin Anton}\input{../tex-inputs/latex-leseansicht-abspann}
